\documentclass{article}
\usepackage{natbib}
\usepackage{hyperref}

\title{Complex QA and Language Models Hybrid Architectures Survey}
\author{Your Name}
\date{\today}

\begin{document}

\maketitle

\begin{abstract}
This paper reviews the state-of-the-art of language model architectures and strategies for complex question-answering (QA, CQA, CPS) with a focus on hybridization. It discusses the limitations of Large Language Models (LLM) in addressing specific and complex questions, and explores recent advancements and challenges in the field.
\end{abstract}

\tableofcontents

\section{Introduction}
Recent advancements in natural language processing (NLP) have led to the development of Large Language Models (LLMs) such as GPT-4 and beyond, capable of understanding and generating human-like text. While these models excel in standard question-answering tasks, they face significant challenges when applied to complex questions that require nuanced understanding across domains and cultures.

Complex question-answering (QA) involves addressing queries that demand deep contextual knowledge, multi-step reasoning, and domain-specific expertise. Examples include understanding cultural variations in concepts like personal freedom or devising optimal strategies for mitigating climate change through a mix of power generation methods. Current LLMs often struggle with these tasks due to limitations in knowledge representation, sensitivity to data biases, and lack of explainability in their decision-making processes.

This survey paper aims to explore the strategies and architectures that enhance LLMs' capabilities in complex QA scenarios, with a particular emphasis on hybrid approaches that integrate symbolic reasoning, structured knowledge bases, and human-in-the-loop feedback mechanisms.

\section{Skills and Evaluation Techniques}
To effectively tackle complex QA tasks, several key skills and evaluation techniques are essential. These include:

\begin{itemize}
    \item \textbf{Domain Adaptation}: Adapting language models to specific domains or cultures to improve accuracy and relevance in answers.
    \item \textbf{Decomposition and Multi-step QA}: Breaking down complex queries into smaller, manageable parts and aggregating results through multi-step reasoning.
    \item \textbf{Safety and Data Protection}: Ensuring models are sensitive to diverse data sources and protect against biases and ethical considerations.
    \item \textbf{Explainability and Truthfulness}: Developing methods to explain model decisions and ensure answers are factually accurate and transparent.
\end{itemize}

These techniques are crucial for evaluating the robustness and reliability of language models in complex QA settings.

\section{Challenges in Complex QA}
Addressing complex QA poses several challenges that go beyond traditional question-answering tasks. Some of the prominent challenges include:

\begin{itemize}
    \item \textbf{Long-form and Non-factoid QA}: Handling queries that require detailed explanations or opinions rather than simple factual answers.
    \item \textbf{Multimodal Search}: Integrating information across multiple modalities such as text, images, and videos to provide comprehensive answers.
    \item \textbf{Safety and Multi-sensitivity Data Protection}: Ensuring models handle sensitive information appropriately and avoid generating harmful or biased outputs.
    \item \textbf{Temporal Reasoning}: Understanding changes and developments over time to provide accurate and up-to-date answers.
\end{itemize}

These challenges necessitate innovative approaches in model design and training strategies.

\section{Current Solutions and Research Trends}
Recent advancements in addressing complex QA include the development of hybrid LLM architectures, integrating techniques such as:

\begin{itemize}
    \item \textbf{Neuro-symbolic and Structured Knowledge Grounding}: Combining neural network approaches with symbolic reasoning to enhance contextual understanding.
    \item \textbf{Program Synthesis}: Automatically generating programs or scripts to solve complex queries requiring procedural knowledge.
    \item \textbf{Iterated Decomposition}: Iteratively breaking down complex questions into simpler sub-problems and aggregating results.
\end{itemize}

These approaches aim to improve model robustness, explainability, and adaptability across diverse QA scenarios.

\section{Conclusion}
In conclusion, the landscape of complex question-answering with language models presents both challenges and opportunities. While current LLMs have demonstrated remarkable capabilities in standard QA tasks, addressing complex queries requires novel architectures and strategies. Hybrid approaches that integrate symbolic reasoning, structured knowledge, and human feedback show promising directions for future research. However, significant challenges such as data sensitivity, explainability, and multi-modal integration remain areas of active investigation.

\section*{Acknowledgments}
We acknowledge the contributions of the research community and the authors of referenced papers in advancing the field of complex QA and language model architectures.

\bibliographystyle{plainnat}
\bibliography{prompt1_try5_35}

\end{document}
